\section{Zadanie 6}
Niech \(X,Y\) będą niezależnymi zmiennymi z rozkładu \(N(0,1)\). Oblicz
\[
	E(X+Y| X > 0, Y>0).
\]

\subsection*{Rozwiązanie}
Z liniowości wartości oczekiwanej i niezależności \(X,Y\):\\
\(E(X+Y|X>0, Y>0) = E(X|X>0, Y>0) + E(Y|X>0, Y>0) = E(X|X>0) + E(Y|Y>0)\)\\
\(X\) i \(Y\) mają ten sam rozkład zatem:\\
\[
	E(X+Y|X>0, Y>0) = 2E(X|X>0)
\]
\[
	E(X|X>0) = \frac{1}{P(X>0)}\int_0^{\infty}{xf(x)dx}
\]
Obliczymy teraz całkę\\
\[
	\int{\frac{1}{\sqrt{2\pi}}xe^{-\frac{1}{2}x^2}dx} = -\frac{1}{\sqrt{2\pi}}e^{-\frac{1}{2}x^2}
\]
\[
	\int_0^{\infty}{\frac{1}{\sqrt{2\pi}}xe^{-\frac{1}{2}x^2}dx} = \frac{1}{\sqrt{2\pi}}
\]
Z symetrii rozkładu normalnego \(P(X>0) = \frac{1}{2}\)\\
Podsumowując:\\
\[
	E(X+Y|X>0, Y>0) = 2 \cdot \frac{1}{\frac{1}{2}}\cdot\frac{1}{\sqrt{2\pi}} = 2\sqrt{\frac{2}{\pi}}
\]

\subsection*{Odpowiedź}
\[
	\boxed{\text{\(2\sqrt{\frac{2}{\pi}}\)}}
\]
